\section{Applications}
The simplicity of the matching is guaranteed once the operators are symmetrized
on the Lorentz indices, and all the traces have been subtracted.
Traceless operators for any $n$ are not difficult to write down.
For example, for $n=4$, we have
\be
\widehat{O}_{4444} = O_{4444} - \frac{3}{4}O_{\left\{\alpha \alpha 4 4 \right\}}
+ \frac{1}{16} O_{\left\{\alpha \alpha \beta \beta \right\}}\,.
\ee
The hadronic matrix element of $\widehat{O}_{4444}(t)$ is calculable
with no external spatial momentum, and the matching with the
$\MS$ scheme at renormalization scale $\mu$ is obtained
using Eq.~\eqref{eq:x_MS}.

The continuum limit of the reduced hadronic matrix element of $\widehat{\rO}_n$,
denoted as $A_n^h(t) = \left\langle x^{n-1} \right\rangle^h(t)$,
would require a determination of $Z_\chi$ from lattice QCD simulations.
However, this can be avoided using the second moment
$A_2^{h,\MS}(\mu) = \left\langle x \right\rangle_{\MS}^h(\mu)$ as an
observable input already computed
with any method of choice. Then all the other moments for $n>2$
are calculable from
\be
\left\langle x^{n-1} \right\rangle_{\MS}^h(\mu) =
\left\langle x^{n-2} \right\rangle_{\MS}^h(\mu)
\frac{c_{n-1}(t,\mu)}{c_n(t,\mu)} R_n^h(t)\,,
\label{eq:ratio_n2}
\ee
where $c_n(t,\mu)$ is given by
Eqs.~(\ref{eq:cn_1l}-\ref{eq:finite}).
The ratios
\be
R_n^h(t) = \frac{\left\langle x^{n-1} \right\rangle^{h}(t)}{\left\langle x^{n-2}\right\rangle^{h}(t)}\,,
\qquad n>2\,,
\label{eq:Rn}
\ee
have a finite continuum limit,
and there is no need to renormalize the flowed fields.
Moreover, these ratios exhibit O($a^2$) scaling violations
once the lattice theory has been nonpertubatively improved.
If convenient numerically, it is also possible to use different hadrons
for the calculation of the moment matrix elements in
Eq.~\eqref{eq:ratio_n2}
because the matching coefficients are independent on the hadron
for which we calculate the matrix element.

The method proposed in this paper can be
also generalized to study distribution amplitudes,
but care should be taken when studying the
short flow time expansion due to the presence of operators
with vanishing contributions for forward matrix
elements~\cite{Bali:2017ude,RQCD:2019osh}.

The gradient flow has previously been employed to propose a locally smeared
operator product expansion in scalar field theory~\cite{Monahan:2015lha}
and to define quasi-PDF in~\cite{Monahan:2016bvm},
with the 1-loop matching calculated in~\cite{Monahan:2017hpu}.
For a more realistic application of the gradient flow
in quasi-PDF calculations see Ref.~\cite{HadStruc:2022yaw}.
While, in principle, the determination of moments of PDFs
should allow a precise reconstruction of the PDFs,
it would be intriguing to explore whether the determination of the
PDF moments outlined in this work
can be combined with quasi-PDF calculations.
